% GENERATED FILE. Edit canonical lessons, then run npm run generate:books.
% canonical-source-hash: fnv1a64:19dc8df3249ee4ec
% canonical-lessons: PA-C256-dakia, PA-C256-parsal, PA-C256-pahunchauna, PA-C256-ghallna, PA-C256-parna

\chapter{The Postman's Round}
\label{ch:pa-a2-the-postman-s-round}
\hlchaptermodality{256}

\begin{chapteropening}
\textbf{By the end of this chapter:} \emph{I can say how letters and parcels reach a house.}

Complete the last lesson of chapter 256: I can say how letters and parcels reach a house.
\end{chapteropening}

\section[\texorpdfstring{\pa{ਡਾਕੀਆ} (ḍākīā)}{ਡਾਕੀਆ (ḍākīā)}]{\pa{ਡਾਕੀਆ} (ḍākīā) — a postman}

\label{lesson:PA-C256-dakia}



\begin{quote}
Before the new one: say the Punjabi for paper, then the Punjabi for to print.
\end{quote}

\subsection*{You'll want to know: \pa{ਡਾਕੀਆ}}
\textbf{\pa{ਡਾਕੀਆ}} — \emph{ḍākīā} — \textquotedblleft{}a postman\textquotedblright{}.

\begin{quote}
\pa{ਡਾਕੀਆ} \pa{ਰੋਜ਼} \pa{ਸਾਡੀ} \pa{ਗਲੀ} \pa{ਵਿੱਚ} \pa{ਆਉਂਦਾ} \pa{ਹੈ।}
\end{quote}

\begin{quote}
\emph{ḍākīā roz sāḍī galī vich āundā hai.}
\end{quote}

\begin{quote}
The postman comes to our lane every day.
\end{quote}

A masculine noun.

\begin{cousinweb}
\textbf{\pa{ਡਾਕੀਆ}} (\emph{ḍākīā}) is built on \emph{ḍāk}, \textquotedblleft{}the post\textquotedblright{}, with the ending \emph{-īā} for the person who does a job; for \emph{ḍāk}, CDIAL gives no Punjabi form (CDIAL 5519).
\end{cousinweb}

\begin{grammarlens}[title={What you've built}]
One more word for reading notices and messages.
\end{grammarlens}

\subsection*{Your turn}
\begin{itemize}
\raggedright
  \item \emph{Say it:} \emph{ḍākīā}
  \item \emph{Say it:} \emph{ḍākīā}, once more
  \item \emph{Say it:} \emph{ḍākīā roz sāḍī galī vich āundā hai.} — The postman comes to our lane every day.
  \item \emph{Recall:} say the Punjabi for paper, then the Punjabi for to print, then say \emph{ḍākīā} again
\end{itemize}

\begin{tcolorbox}[breakable,colback=teal!4,colframe=teal!35!black,title={Before you move on}]
What does \pa{ਡਾਕੀਆ} mean? (\textquotedblleft{}A postman\textquotedblright{}.) Say it once more.
\end{tcolorbox}
\section[\texorpdfstring{\pa{ਪਾਰਸਲ} (pārsal)}{ਪਾਰਸਲ (pārsal)}]{\pa{ਪਾਰਸਲ} (pārsal) — a parcel}

\label{lesson:PA-C256-parsal}



\begin{quote}
Before the new one: say the Punjabi for to print, then the Punjabi for a postman.
\end{quote}

\subsection*{You'll want to know: \pa{ਪਾਰਸਲ}}
\textbf{\pa{ਪਾਰਸਲ}} — \emph{pārsal} — \textquotedblleft{}a parcel\textquotedblright{}.

\begin{quote}
\pa{ਅਮਨ} \pa{ਨੇ} \pa{ਮੈਨੂੰ} \pa{ਇੱਕ} \pa{ਪਾਰਸਲ} \pa{ਭੇਜਿਆ।}
\end{quote}

\begin{quote}
\emph{Aman ne mainū̃ ikk pārsal bhejiā.}
\end{quote}

\begin{quote}
Aman sent me a parcel.
\end{quote}

A masculine noun.

\begin{cousinweb}
\textbf{\pa{ਪਾਰਸਲ}} (\emph{pārsal}) is borrowed from English \emph{parcel}.
\end{cousinweb}

\begin{grammarlens}[title={What you've built}]
One more word for reading notices and messages.
\end{grammarlens}

\subsection*{Your turn}
\begin{itemize}
\raggedright
  \item \emph{Say it:} \emph{pārsal}
  \item \emph{Say it:} \emph{pārsal}, once more
  \item \emph{Say it:} \emph{Aman ne mainū̃ ikk pārsal bhejiā.} — Aman sent me a parcel.
  \item \emph{Recall:} say the Punjabi for to print, then the Punjabi for a postman, then say \emph{pārsal} again
  \item \emph{Write it:} copy \textbf{\pa{ਪਾਰਸਲ}} once, checking the whole word
\end{itemize}

\begin{tcolorbox}[breakable,colback=teal!4,colframe=teal!35!black,title={Before you move on}]
What does \pa{ਪਾਰਸਲ} mean? (\textquotedblleft{}A parcel\textquotedblright{}.) Say it once more.
\end{tcolorbox}
\section[\texorpdfstring{\pa{ਪਹੁੰਚਾਉਣਾ} (pahunchāuṇā)}{ਪਹੁੰਚਾਉਣਾ (pahunchāuṇā)}]{\pa{ਪਹੁੰਚਾਉਣਾ} (pahunchāuṇā) — to deliver; to take something to a place}

\label{lesson:PA-C256-pahunchauna}



\begin{quote}
Before the new one: say the Punjabi for a postman, then the Punjabi for a parcel.
\end{quote}

\subsection*{You'll want to know: \pa{ਪਹੁੰਚਾਉਣਾ}}
\textbf{\pa{ਪਹੁੰਚਾਉਣਾ}} — \emph{pahunchāuṇā} — \textquotedblleft{}to deliver; to take something to a place\textquotedblright{}.

\begin{quote}
\pa{ਡਾਕੀਆ} \pa{ਚਿੱਠੀਆਂ} \pa{ਸਾਡੇ} \pa{ਘਰ} \pa{ਪਹੁੰਚਾਉਂਦਾ} \pa{ਹੈ।}
\end{quote}

\begin{quote}
\emph{ḍākīā chiṭṭhīāṁ sāḍe ghar pahunchāundā hai.}
\end{quote}

\begin{quote}
The postman delivers the letters to our house.
\end{quote}

Past (I …): \pa{ਮੈਂ} \pa{ਪਾਰਸਲ} \pa{ਪਹੁੰਚਾਇਆ}, \emph{maiṁ pārsal pahunchāiā}.

\begin{cousinweb}
\textbf{\pa{ਪਹੁੰਚਾਉਣਾ}} (\emph{pahunchāuṇā}) is the causative partner of \emph{pahũcṇā}, \textquotedblleft{}to arrive\textquotedblright{}; CDIAL lists Punjabi \emph{pahũcṇā} under Sanskrit \emph{prabhūta}, \textquotedblleft{}much, great\textquotedblright{} (CDIAL 8716).
\end{cousinweb}

\begin{grammarlens}[title={What you've built}]
One more word for reading notices and messages.
\end{grammarlens}

\subsection*{Your turn}
\begin{itemize}
\raggedright
  \item \emph{Say it:} \emph{pahunchāuṇā}
  \item \emph{Say it:} \emph{pahunchāuṇā}, once more
  \item \emph{Say it:} \emph{ḍākīā chiṭṭhīāṁ sāḍe ghar pahunchāundā hai.} — The postman delivers the letters to our house.
  \item \emph{Recall:} say the Punjabi for a postman, then the Punjabi for a parcel, then say \emph{pahunchāuṇā} again
  \item \emph{Write it:} copy \textbf{\pa{ਪਹੁੰਚਾਉਣਾ}} once, checking the whole word
\end{itemize}

\begin{tcolorbox}[breakable,colback=teal!4,colframe=teal!35!black,title={Before you move on}]
What does \pa{ਪਹੁੰਚਾਉਣਾ} mean? (\textquotedblleft{}To deliver; to take something to a place\textquotedblright{}.) Say it once more.
\end{tcolorbox}
\section[\texorpdfstring{\pa{ਘੱਲਣਾ} (ghallṇā)}{ਘੱਲਣਾ (ghallṇā)}]{\pa{ਘੱਲਣਾ} (ghallṇā) — to send off; to send}

\label{lesson:PA-C256-ghallna}



\begin{quote}
Before the new one: say the Punjabi for a parcel, then the Punjabi for to deliver.
\end{quote}

\subsection*{You'll want to know: \pa{ਘੱਲਣਾ}}
\textbf{\pa{ਘੱਲਣਾ}} — \emph{ghallṇā} — \textquotedblleft{}to send off; to send\textquotedblright{}.

\begin{quote}
\pa{ਮੈਂ} \pa{ਹਰ} \pa{ਮਹੀਨੇ} \pa{ਪਿੰਡ} \pa{ਚਿੱਠੀ} \pa{ਘੱਲਦਾ} \pa{ਹਾਂ।}
\end{quote}

\begin{quote}
\emph{maiṁ har mahīne piṇḍ chiṭṭhī ghalldā hāṁ.}
\end{quote}

\begin{quote}
Every month I send a letter to the village.
\end{quote}

Past (I …): \pa{ਮੈਂ} \pa{ਅਮਨ} \pa{ਨੂੰ} \pa{ਪਾਰਸਲ} \pa{ਘੱਲਿਆ}, \emph{maiṁ Aman nū̃ pārsal ghalliā}.

\begin{cousinweb}
\textbf{\pa{ਘੱਲਣਾ}} (\emph{ghallṇā}) is inherited from a reconstructed \emph{*ghalyati}, \textquotedblleft{}pours\textquotedblright{}; CDIAL lists Punjabi \emph{ghallṇā}, \textquotedblleft{}to send\textquotedblright{}, under it (CDIAL 4457).
\end{cousinweb}

\begin{grammarlens}[title={What you've built}]
One more word for reading notices and messages.
\end{grammarlens}

\subsection*{Your turn}
\begin{itemize}
\raggedright
  \item \emph{Say it:} \emph{ghallṇā}
  \item \emph{Say it:} \emph{ghallṇā}, once more
  \item \emph{Say it:} \emph{maiṁ har mahīne piṇḍ chiṭṭhī ghalldā hāṁ.} — Every month I send a letter to the village.
  \item \emph{Recall:} say the Punjabi for a parcel, then the Punjabi for to deliver, then say \emph{ghallṇā} again
  \item \emph{Write it:} copy \textbf{\pa{ਘੱਲਣਾ}} once, checking the whole word
\end{itemize}

\begin{tcolorbox}[breakable,colback=teal!4,colframe=teal!35!black,title={Before you move on}]
What does \pa{ਘੱਲਣਾ} mean? (\textquotedblleft{}To send off; to send\textquotedblright{}.) Say it once more.
\end{tcolorbox}
\section[\texorpdfstring{\pa{ਪਾੜਨਾ} (pāṛnā)}{ਪਾੜਨਾ (pāṛnā)}]{\pa{ਪਾੜਨਾ} (pāṛnā) — to tear}

\label{lesson:PA-C256-parna}



\begin{quote}
Before the new one: say the Punjabi for to deliver, then the Punjabi for to send off.
\end{quote}

\subsection*{You'll want to know: \pa{ਪਾੜਨਾ}}
\textbf{\pa{ਪਾੜਨਾ}} — \emph{pāṛnā} — \textquotedblleft{}to tear\textquotedblright{}.

\begin{quote}
\pa{ਲਿਫ਼ਾਫ਼ਾ} \pa{ਖੋਲ੍ਹੋ} \pa{ਪਰ} \pa{ਚਿੱਠੀ} \pa{ਨਾ} \pa{ਪਾੜੋ।}
\end{quote}

\begin{quote}
\emph{lifāfā kholho par chiṭṭhī nā pāṛo.}
\end{quote}

\begin{quote}
Open the envelope, but do not tear the letter.
\end{quote}

Past (I …): \pa{ਮੈਂ} \pa{ਪੁਰਾਣਾ} \pa{ਕਾਗ਼ਜ਼} \pa{ਪਾੜਿਆ}, \emph{maiṁ purāṇā kāghaz pāṛiā}.

\begin{cousinweb}
\textbf{\pa{ਪਾੜਨਾ}} (\emph{pāṛnā}) is inherited from Sanskrit \emph{pāṭayati}, \textquotedblleft{}splits\textquotedblright{}; CDIAL lists Punjabi \emph{pāṛnā}, \textquotedblleft{}to split, tear\textquotedblright{}, under it (CDIAL 8032).
\end{cousinweb}

\begin{grammarlens}[title={What you've built}]
One more word for reading notices and messages.
\end{grammarlens}

\subsection*{Your turn}
\begin{itemize}
\raggedright
  \item \emph{Say it:} \emph{pāṛnā}
  \item \emph{Say it:} \emph{pāṛnā}, once more
  \item \emph{Say it:} \emph{lifāfā kholho par chiṭṭhī nā pāṛo.} — Open the envelope, but do not tear the letter.
  \item \emph{Recall:} say the Punjabi for to deliver, then the Punjabi for to send off, then say \emph{pāṛnā} again
  \item \emph{Write it:} copy \textbf{\pa{ਪਾੜਨਾ}} once, checking the whole word
\end{itemize}

\begin{tcolorbox}[breakable,colback=teal!4,colframe=teal!35!black,title={Before you move on}]
What does \pa{ਪਾੜਨਾ} mean? (\textquotedblleft{}To tear\textquotedblright{}.) Say it once more.
\end{tcolorbox}
